% Étape 3 (1/2) : le bloc d'attention, en étapes (FreeFlow, éq. 3)
% Compiler avec pdflatex (ou Overleaf).
\documentclass[border=10pt,varwidth=700pt]{standalone}
\usepackage[T1]{fontenc}
\usepackage[utf8]{inputenc}
\usepackage[scaled]{helvet}
\renewcommand{\familydefault}{\sfdefault}
\usepackage{sansmath}
\sansmath
\usepackage{amsmath}
\usepackage{array}
\usepackage{xcolor}
\usepackage{colortbl}

\definecolor{orange}{HTML}{FFD0BF}
\definecolor{creme}{HTML}{FFF6EC}

\newcommand{\num}[1]{\textbf{#1}}
\renewcommand{\arraystretch}{2.1}
\arrayrulecolor{orange}

\begin{document}
\fontsize{13}{16}\selectfont
\begin{tabular}{>{\centering\arraybackslash}m{22pt} >{\raggedright\arraybackslash}m{95pt} >{\raggedright\arraybackslash}m{215pt} >{\centering\arraybackslash}m{62pt} >{\raggedright\arraybackslash}m{250pt}}
\rowcolor{creme}
 & \textbf{Étape} & \textbf{Formule} & \textbf{Taille} & \textbf{Ce que ça fait} \\
\hline
\num{1} & Normaliser & $\tilde{X} = \mathrm{Norm}(X_w)$ & $n \times D$ & met les valeurs à la même échelle \\
\hline
\num{2} & Projeter & $Q = \tilde{X} W_Q,\ \ K = \tilde{X} W_K,\ \ V = \tilde{X} W_V$ & $n \times D$ & $Q$ : ce que le patch cherche \newline $K$ : son étiquette \quad $V$ : son contenu \\
\hline
\num{3} & Scores & $S = Q K^{\top}$, \ \ $S[j,k] = q_j \cdot k_k$ & $n \times n$ & ressemblance entre les patchs $j$ et $k$ \\
\hline
\num{4} & Poids & $A = \mathrm{Softmax}\!\left(\dfrac{\log n}{\sqrt{64}}\, S\right)$ & $n \times n$ & chaque ligne devient des poids (somme = 1) \\
\hline
\num{5} & Mélange & $O = A\, V$, \ \ $o_j = \sum_k A[j,k]\, v_k$ & $n \times D$ & chaque patch récupère le contenu des patchs \newline qu'il regarde \\
\hline
\end{tabular}

\vspace{10pt}
\begin{minipage}{660pt}
\fontsize{12}{15}\selectfont
$n = N/16$ (Win, Swin) ou $N/4$ (Global) \qquad $W_Q, W_K, W_V$ : $D \times D$, appris \qquad $\sqrt{64}$ : taille d'une tête \\ $\log n$ : garde le Softmax net quand $n$ est grand

\vspace{4pt}
Entrée $X_w$ : $n \times D$ \ $\rightarrow$ \ Sortie $O$ : $n \times D$ (même taille) \ $\rightarrow$ têtes, résiduel et MLP (2/2)
\end{minipage}
\end{document}
